\documentclass[tikz,border=8pt]{standalone}
\usepackage{amsmath}
\usetikzlibrary{calc}

\begin{document}
\begin{tikzpicture}[font=\small]

\def\globalscale{0.18}
\def\labely{-4.5} % �� 모든 label의 공통 y좌표

\newcommand{\CircleExample}[3]{%
\begin{scope}[shift={(#3,0)}]

    \pgfmathsetmacro{\b}{#1}
    \pgfmathsetmacro{\d}{#2}
    \pgfmathsetmacro{\R}{\globalscale*\d}

    \pgfmathsetmacro{\theta}{2*asin(\b/\d)}
    \pgfmathsetmacro{\ratio}{360/\theta}
    \pgfmathtruncatemacro{\N}{ceil(\ratio)}
    \pgfmathtruncatemacro{\NmOne}{\N-1}

    \pgfmathsetmacro{\ptsize}{ifthenelse(\N>120,0.45,
                             ifthenelse(\N>50,0.65,
                             ifthenelse(\N>20,0.9,1.5)))}

    % circle
    \draw[gray!70] (0,0) circle (\R);

    % points
    \foreach \k in {0,...,\NmOne} {
        \fill ({\R*cos(90 + \k*\theta)},
              {\R*sin(90 + \k*\theta)})
              circle (\ptsize pt);
    }

    % �� 고정된 y 위치로 label
    \node[font=\Large] at (0,\labely) {$\displaystyle (b,d)=(#1,#2)$};

\end{scope}
}

\CircleExample{1}{10}{0}
\CircleExample{1}{8}{4.8}
\CircleExample{9.5}{11}{9.5}
\CircleExample{3}{15}{15.5}

\end{tikzpicture}
\end{document}